% ECONOMICS_PROBLEM: {"id": "C24-509", "title": "Nonrandom attrition and missing-not-at-random outcomes", "jel_code": "C24", "source_id": "OP-0109"}
% !TeX program = xelatex
\documentclass[11pt,a4paper]{article}
\usepackage[margin=23mm]{geometry}
\usepackage{fontspec}
\setmainfont{Latin Modern Roman}
\usepackage{amsmath,amssymb,mathtools}
\usepackage{xcolor,enumitem,booktabs,longtable,array,xurl,hyperref}
\definecolor{muted}{HTML}{53636C}
\hypersetup{colorlinks=true,urlcolor=blue,linkcolor=blue}
\setlength{\parindent}{0pt}
\setlength{\parskip}{5pt}
\newcommand{\R}{\mathbb R}
\newcommand{\E}{\mathbb E}
\newcommand{\N}{\mathbb N}
\newcommand{\ind}{\mathbf 1}
\newcommand{\Var}{\operatorname{Var}}
\newcommand{\Cov}{\operatorname{Cov}}
\newcommand{\supp}{\operatorname{supp}}
\DeclareMathOperator*{\argmin}{arg\,min}
\DeclareMathOperator*{\argmax}{arg\,max}
\newcommand{\entrymeta}[3]{{\sffamily\small\color{muted}\textbf{\texttt{#1}}\quad|\quad #2\par #3\par}}
\newcommand{\Problemref}[1]{\texttt{#1}}
\newcommand{\JELref}[2]{\texttt{#2} (JEL \texttt{#1})}
\begin{document}
\section*{Nonrandom attrition and missing-not-at-random outcomes}
\textbf{Problem ID:} \texttt{C24-509}\quad \textbf{JEL:} \texttt{C24}\par
% BEGIN ECONOMICS STATEMENT
\label{problem:OP-0109}
\label{atlas:prob:E9}

\noindent\textbf{Original identifier:} E9 (atlas item 109). \textbf{Field:} Econometrics and causal inference. \textbf{Classification:} editorial research family with established restricted solutions; current open status not certified.

\subsection*{Definitions and assumptions}

Let an independent randomized trial observe $(D,X,R,RY)$, where $D\in\{0,1\}$ has known propensity in $[\eta,1-\eta]$ for $\eta>0$, $R\in\{0,1\}$ indicates outcome observation, and $Y=Y(D)$ when observed. Assume unconditional random assignment for the unadjusted formulas below. Potential outcomes satisfy $Y(d)\in[a,b]$, with finite known $a<b$, and potential observation indicators are $R(d)$. Define population ATE $\theta=\mathbb E[Y(1)-Y(0)]$. Write $p_d=\Pr(R=1\mid D=d)$ and $\mu_d=\mathbb E[Y\mid D=d,R=1]$ when $p_d>0$. The observed conditional mean is not the full-population potential-outcome mean.

\subsection*{Formal research question}

For assumptions indexed by a sensitivity parameter $\delta\in[0,1]$, determine sharp bounds and honest inference for both population and principal-stratum targets. With randomization and only bounded outcomes, define $L_d=p_d\mu_d+(1-p_d)a$ and $U_d=p_d\mu_d+(1-p_d)b$, with the observed component taken as zero when $p_d=0$. Then
\[
\theta\in[L_1-U_0,\ U_1-L_0].
\]
Under selection monotonicity $R(1)\ge R(0)$ almost surely, define $\theta_A=\mathbb E[Y(1)-Y(0)\mid R(1)=R(0)=1]$, provided that this stratum has positive probability. Derive trimming bounds for $\theta_A$ and their extension when $\Pr\{R(1)<R(0)\}\le\delta$. Construct simultaneous confidence regions uniform over selection rates approaching zero and distributions with mass points.

\subsection*{Interpretation and scope}

Monotone selection does not say outcomes increase with treatment. Lee-type trimming identifies bounds for an always-observed principal stratum; relabeling them as bounds for the population ATE is incorrect. Missing-at-random permits weighting only under conditional independence and positivity; a predictive imputation model cannot make these assumptions empirically true. For wages observed only in employment, an all-population wage outcome may be undefined, so a composite outcome or principal-stratum target must be declared. Refreshment samples or randomized follow-up can add information, but their own observation mechanisms require specification. An honest report distinguishes identification width from sampling uncertainty and displays conclusions as assumptions vary.

\subsection*{Source audit and status}

[S1] studies selective observation, [S2] derives sharp trimming bounds under monotonicity, and [S3] is the verified second-edition missing-data monograph. Existing nonparametric and model-based methods already address restricted cases. The sensitivity-and-uniform-inference extension is an editorial target, not a claim that the original generic attrition problem lacks established solutions or remains universally open in 2026.

\subsection*{References}

\begin{enumerate}[label={[\textnormal{S}\arabic*]},leftmargin=*]

\item Charles F. Manski (1989). \emph{Anatomy of the Selection Problem}. Journal of Human Resources 24(3), 343–360. \url{https://doi.org/10.2307/145818}. \textbf{Locator:} Journal issue table of contents, lines 26–34. \textbf{Support and limits:} Selective outcome observation and nonparametric identification; journal issue record verifies title/author/year.

\item David S. Lee (2009). \emph{Training, Wages, and Sample Selection: Estimating Sharp Bounds on Treatment Effects}. Review of Economic Studies 76(3), 1071–1102. \url{https://doi.org/10.1111/j.1467-937X.2009.00536.x}. \textbf{Locator:} Publisher or institutional abstract and bibliographic record. \textbf{Support and limits:} Sharp trimming bounds under selection monotonicity; target is always-observed principal stratum, not general population ATE.

\item Roderick J. A. Little and Donald B. Rubin (2002). \emph{Statistical Analysis with Missing Data}. Wiley, second edition. \url{https://doi.org/10.1002/9781119013563}. \textbf{Locator:} Publisher or institutional abstract and bibliographic record. \textbf{Support and limits:} Missing-data theory including ignorable/nonignorable mechanisms; MNAR is not identified from observed data alone.

\end{enumerate}

\subsection*{Common conventions: Source mathematical conventions}

Unless an entry states otherwise, agent and firm sets are finite. State and action spaces are measurable, strategies and outcome maps are measurable, probability laws are specified, and expectations exist for the targets under discussion. Infinite-horizon discount factors lie in $(0,1)$ and discounted objectives are finite. Norms, tolerances and complexity input sizes must be specified for computational comparisons. Constants or primitives that appear locally are inputs, not empirical facts asserted by this catalogue.

\textbf{Identification.} A statistical model $m\in\mathcal M$ implies an observable law $P_m$ and target $\theta(m)$. At law $P$, its identified set is
\[
\Theta(P;\mathcal M)=\{\theta(m):m\in\mathcal M,\ P_m=P\}.
\]
Point identification holds when this set is a singleton, or equivalently when $P_{m_1}=P_{m_2}$ implies $\theta(m_1)=\theta(m_2)$ for all compatible models. Sharp bounds mean bounds attained, or approached when the boundary is not attained, by compatible models. Writing this set is a definition, not a solution: the substantive task is to establish informative restrictions and derive or estimate the set. An unrestricted-confounding model may give only trivial bounds.

\textbf{Inference.} Identification, estimability and valid uncertainty quantification are separate questions. Fix $\alpha\in(0,1)$. A confidence procedure $C_n$ over a declared law class $\mathcal P$ has asymptotic uniform coverage at level $1-\alpha$ when
\[
\liminf_{n\to\infty}\inf_{P\in\mathcal P}\Pr_P\{\theta(P)\in C_n\}\geq1-\alpha.
\]
For set-identified targets the coverage event must be defined separately, for example coverage of the full identified set. If an entry uses identification-indexed models instead of uniquely indexed laws, the same coverage convention is applied over those models. Pointwise limits do not establish uniform validity, and asymptotic validity does not guarantee finite-sample precision. Sample sizes, dependence structures and weak-identification sequences are part of each application.

\textbf{Counterfactuals.} $Y_i(a)$ is a potential outcome under a specified intervention, including its timing, implementation and versions. It is not $Y_i$ conditional on voluntarily choosing $a$. A full intervention vector permits interference; shorthand scalar treatment notation does not silently impose no interference. Assignment, selection, attrition, anticipation and measurement must be specified. A claim about an instrument's local effect is not automatically a population policy effect.

\textbf{Economic equilibrium.} A counterfactual model includes best responses, constraints, feasibility, market clearing, state transitions and expectations consistent with the stated information. When several equilibria exist, either a declared selector is an input or the target is set-valued. Comparisons across policy regimes must use the stated selection convention. Reduced-form relationships are not presumed invariant after an equilibrium-changing intervention.

\textbf{Optimization and welfare.} Optimization is over the stated feasible set. If attainment is not established, $\sup$ or $\inf$ and $\varepsilon$-optimal choices replace an asserted maximizer or minimizer. For example, with value $V=\sup_{a\in\mathcal A}W(a)<\infty$, an $\varepsilon$-optimal set is $\{a\in\mathcal A:W(a)\geq V-\varepsilon\}$ for $\varepsilon>0$. Benefits, costs, risk and health or environmental outcomes can be aggregated only using declared units and conversion weights. Interpersonal utility weights, the treatment of fiscal transfers and foreign profits, discounting, and the behavioral welfare criterion are explicit inputs. A comparative-static derivative additionally requires existence and appropriate differentiability of the selected counterfactual.

\textbf{Sources.} Local labels [S1], [S2], and so forth restart in every question. Locators identify the relevant abstract, model, section or result at the level actually checked; a bibliographic or abstract-level verification is not represented as inspection of an entire proof. Working papers are distinguished from later publications and changing versions. Source summaries are paraphrased. All URL checks and editorial formulations refer to the audit date stated above.
% END ECONOMICS STATEMENT
\end{document}
